% 本文件由 LaTeX 排版 Agent 根据有效 translation.json 自动生成。
% author=Tertullian; work_id=0302; title=论偶像崇拜

\chapter{论偶像崇拜}

% structure: p0001 source=h1 target=omitted reason=page-title-already-rendered-by-parent

% structure: p0002 source=h2 target=section reason=relative-heading-rank; base=h2

\section{第一章 偶像崇拜一词的广泛范围}

人类的首要罪行，加于世界的最重罪责，导致审判的全部起因，是偶像崇拜。因为，虽然每一单个的过犯都保有自己特有的性质，虽然它也将按自己特有的名目受审，然而它仍被归在偶像崇拜这\Emph{一般}账目之下。放下名目，查看行为，拜偶像者同样是个杀人犯。你问他是杀了谁？如果这对加重控诉有丝毫贡献，那么他杀的不是陌生人，也不是私敌，而是他自己。用什么陷阱？他错误的陷阱。用什么凶器？对天主的冒犯。打多少下？就像他的偶像崇拜一样多。凡断言拜偶像者\Emph{不灭亡}的人，也将断言拜偶像者没有犯杀人罪。此外，你可以在同一罪行中认出\Emph{奸淫}和\Emph{淫乱}；因为侍奉假神的人无疑是真理的\Emph{奸淫者}，因为一切虚假都是奸淫。同样，他也沉溺于淫乱。因为谁与邪灵同工，却不徘徊在普遍的污秽和淫乱中呢？正因如此，圣经在责备偶像崇拜时使用淫乱的称呼。我认为\Emph{欺诈}的本质，就是攫取他人之物，或拒绝给予他人所应得的；当然，对\Emph{人}行欺诈是最严重罪行的名号。然而，偶像崇拜对天主行欺诈，拒绝将属于祂的尊荣归于祂，反而归于他人；因此它还在欺诈之外加上\Emph{侮辱}。但如果欺诈与淫乱和奸淫一样招致死亡，那么在这些情况下，偶像崇拜同样无法免于杀人罪的控告。在如此恶毒、如此吞噬救恩的罪行之后，所有其他罪行，在某种程度上，各按其序，也在偶像崇拜中找到了自身的本质。其中也有世界的\Emph{贪欲}。因为有什么偶像崇拜的庆典是没有衣着和装饰的呢？其中有\Emph{淫荡}和\Emph{酗酒}；因为这些庆典多半是为了食物、肚腹和食欲而举行的。其中有\Emph{不义}。因为有什么比它更不义呢？它竟不认识公义之父。其中也有\Emph{虚妄}，因为它的整个体系都是虚妄的。其中有\Emph{谎谬}，因为它的全部实质都是虚假的。因此，所有罪行都在偶像崇拜中被发现，而偶像崇拜也在所有罪行中被发现。即使从另一面看，既然一切过犯都带有反对天主的味道，而凡带有反对天主味道的，没有一样不被归于魔鬼和邪灵，偶像正是它们的财产；那么无疑，凡犯罪的人都要被控以偶像崇拜，因为他做了属于偶像所有者的事。\par

% structure: p0004 source=h2 target=section reason=relative-heading-rank; base=h2

\section{第二章 更狭义上的偶像崇拜 其丰富性}

但是，让罪行的普遍名称退回到它们自己行为的特殊性中去；让偶像崇拜停留在它自身之中。一个与天主为敌的名称，一种如此丰富的罪行本质，它伸出如此多的分枝，散布如此多的脉络，已足以自足；因为从这个\Emph{名称}，大部分情况下，提取了我们必须预防的偶像崇拜之广泛性的所有模式的材料，因为它以多种方式颠覆天主的仆人；这不仅在未被察觉时，而且在被掩盖时也是如此。大多数人简单地认为偶像崇拜只能在这些意义上解释，即：若有人烧香，或献上\Emph{祭牲}，或举行祭宴，或受缚于某些神圣职务或司祭职；就好像一个人认为奸淫只算作亲吻、拥抱和实际肉体接触；或杀人只算作流血和实际夺命。但我们确知主对那些罪行的界定何等广泛：当祂定义奸淫\Emph{亦在于}贪欲，\BibleRef{玛 5:28}\Quote{若有人以淫念注目，}并以不洁的冲动搅动自己的灵魂；当祂判定杀人\BibleRef{玛 5:22}\Emph{亦在于}一句咒骂或斥责的话，以及每一个愤怒的冲动，并忽视对弟兄的爱德，正如若望教导的，\BibleRef{若一 3:15}凡恨弟兄的就是杀人者。否则，魔鬼的恶意狡诈，以及主天主的训导（祂藉此坚固我们对抗魔鬼的深奥，\BibleRef{默 2:24}）都将范围有限，如果我们只按甚至连外邦民族也裁定为可惩罚的过犯受审。我们的义德如何能超过经师和法利塞人的义德，正如主所吩咐的，\BibleRef{玛 5:20}\Quote{义德超过经师和法利塞人的义德，}除非我们看透了那敌对性质（即不义）的丰盛？但若不义之首是偶像崇拜，那么首要之点就是我们预先防备偶像崇拜的丰盛，同时不仅在其明显的表现中认出它。\par

% structure: p0006 source=h2 target=section reason=relative-heading-rank; base=h2

\section{第三章 偶像崇拜：名称的起源与意义}

偶像在古代并不存在。在这怪物的工匠们涌现之前，圣殿独自矗立，神龛空置，正如直到今日在某些地方，古代习俗的痕迹依然留存。然而，偶像崇拜当时已在实行，并非以那名称，而是以那功能；因为即使在今日，它也可以在圣殿之外，且在没有偶像的情况下实行。但当魔鬼将雕像和肖像以及各样形像的工匠引入世界时，那以前粗陋的人类灾难事业便从偶像那里获得了名称和发展。从此，任何以某种方式产生偶像的技艺，立刻成为偶像崇拜的泉源。因为无论模铸工铸造、雕刻工雕刻、还是刺绣工编织\Emph{那偶像}，都无关紧要；因为问题也不在于材料，无论偶像是由石膏、颜料、石头、铜、银还是线制成。因为即使没有偶像，偶像崇拜也已被犯下；当偶像存在时，它是何种样式、何种材料、何种形状，都无关紧要；免得有人认为只有奉为神圣的人形之物才算作偶像。要确立这一点，需要解释这个词。在希腊语中，\OriginalTerm{形式}{Eidos}意指形式；\OriginalTerm{小形式}{Eidolon}是它的指小词，在我们的语言中通过类似过程产生\Quote{小形式}。因此，每一个形式或小形式都当被称为偶像。因此，偶像崇拜就是\Quote{对每一偶像的一切侍奉和服务}。于是，每一位偶像工匠都犯有同一种罪行，除非那为自己祝圣了牛犊的肖像（而不是人的肖像）的百姓，并没有犯下偶像崇拜之罪。\par

% structure: p0008 source=h2 target=section reason=relative-heading-rank; base=h2

\section{偶像不可制造，更不可崇拜。偶像和偶像制造者同属一类}

天主禁止制造偶像，同样禁止崇拜偶像。既然制造可受崇拜之物是先行行为，那么禁止制造（若崇拜为非法）便是先行禁令。为此原因——即根除偶像崇拜的原料——神圣法律宣告：\BibleQuote{你不可制造偶像；}并连同宣告：\BibleQuote{也不制造天上、地下、水中之物的任何肖像，}从而禁止天主的仆人在普天之下从事这类行为。\Person{哈诺客}早已预言：\Quote{魔鬼和背教天使的幽灵，将把一切元素、宇宙的一切装饰、天上、海里、地上的一切事物都转为偶像崇拜，使它们被祝圣为神，与天主对立。}因此，人的错误崇拜一切事物，唯独不崇拜万有的造物主本身。那些事物的肖像是偶像；肖像的祝圣是偶像崇拜。偶像崇拜所招致的任何罪过，必然归咎于每一个偶像的制造者。简言之，同一位\Person{哈诺客}以普遍的威胁同时谴责了崇拜偶像者和制造偶像者。他又说：\Quote{我向你们罪人起誓，为流血之毁灭日，补赎正在预备。你们这些事奉石头、制造金、银、木、石、泥像的人，以及崇拜幽灵、魔鬼和庙宇中的邪灵、以及一切不合知识的错误者，从它们那里得不到任何帮助。}但\Person{依撒意亚}说：\BibleQuote{你们见证，除了我以外，还有神吗？}\BibleQuote{那时，塑造和雕刻它们的人并不存在：全是虚妄！他们做自己喜欢的事，但对他们毫无益处！}而随后整篇讲论既禁止制造者也禁止崇拜者：其结尾是：\BibleQuote{要知道他们的心是灰土，没有人能解救自己的灵魂。}在这句话中，\Person{达味}同样把制造者包括在内。他说：\BibleQuote{愿制造它们的人变成这样。}为什么我这个记忆有限的人还要进一步提议？为什么还要从圣经中回顾更多？仿佛圣神的声音不够，或者还需要进一步思考，主是否优先咒诅和谴责那些祂所咒诅和谴责的崇拜者所制之物的\Emph{制造者}和\Emph{崇拜者}！\par

% structure: p0010 source=h2 target=section reason=relative-heading-rank; base=h2

\section{处理各种异议或借口}

我们当然要更加费心回答这类工匠的托词，他们若对纪律有所认识，本绝不应被接纳进入天主的殿。首先，那句惯常被用来顶撞我们的话——\Quote{此外我无以谋生}——可以更严厉地反驳回去：\Quote{那么，你就有东西赖以生存了吗？若凭你自己的法则，你与天主有何相干？}其次，关于他们竟敢从圣经中引用的论据，即\Quote{宗徒说过：\InnerQuote{各人怎样被找到，就该怎样坚持。}}因此，按那种解释，我们所有人都可以在\Emph{罪}中坚持了！因为我们中没有谁不是作为\Emph{罪人}被找到的，因为基督降世的原因无非就是释放\Emph{罪人}。再者，他们说同一宗徒留下了一条诫命，要他以身作则，\Quote{各人要亲手做工以谋生计}。如果这条诫命适用于\Emph{所有}手，那么我相信连浴场窃贼也靠手维生，强盗也靠手获得生计；伪造者同样用他们的手（当然不是脚，而是手）进行邪恶的书写；而演员不仅靠手，还靠他们的整个肢体来谋生。因此，教会应当向所有靠双手和自己的劳动维生的人敞开，如果那些不被天主的纪律所接纳的技艺没有例外的话。但有人反对我们提出的\Quote{禁止肖像}的主张，说：\Quote{那么，摩西在旷野中为什么用铜造了一条蛇的像？}那些曾作为未来某种隐秘安排的根基形象的像，并非意图废除法律，而是作为其终极目标的预象，属于另一类别。否则，如果我们像法律的敌人们那样解释这些事，我们岂不是也像马尔基翁派那样，将反复无常归诸于全能者？他们正是以这种方式摧毁祂，认为祂是可变的，因为祂在一处禁止，在另一处却命令。但若有人假装不知道，那铜蛇的像，如同被悬挂者一般，预示了主的十字架的形状，那十字架要释放我们脱离蛇——即魔鬼的使者——同时藉着它本身悬挂起被击溃的魔鬼；或无论那形象还有何种其他解释已被启示给更有德之人，这都无关紧要，只要我们记得宗徒确认当时所有事情\Emph{以预像的方式}发生在百姓身上。这就够了：同一位天主，既藉着法律禁止制造肖像，又藉着关于铜蛇的特别诫命禁止肖像。如果你敬畏同一位天主，你就拥有祂的法律：\Quote{你不可制造任何肖像。}如果你也回顾那后来命令制造肖像的诫命，那么你也效法摩西：不要违反法律制造任何像，除非天主也命令了你。\par

% structure: p0012 source=h2 target=section reason=relative-heading-rank; base=h2

\section{第六章 偶像崇拜被圣洗所谴责。制造偶像其实就是崇拜它}

即使没有天主的法律禁止我们制造偶像，即使圣神的声音没有发出普遍的威胁，既针对偶像的制造者，也针对崇拜者，我们也能从我们的圣事本身推导出这类技艺与信仰相悖的结论。因为我们如何\Emph{弃绝}了魔鬼和他的天使，如果我们\Emph{制造}他们？我们与他们宣告了何种决裂？我不说\Emph{与}他们同居，而是\Emph{依赖}他们生活？我们与他们进入了何种不和，而我们还要为生计欠他们的情？你能用舌头否认你用手所承认的吗？用言语废除你用行为所制造的吗？你制造如此众多的神，却宣讲独一的天主？你制造虚假的神，却宣讲真实的天主？有人说：\Quote{我\Emph{制造}，但我不\Emph{崇拜}。}仿佛有一种理由使他不敢\Emph{崇拜}，而不同于他也不应\Emph{制造}的理由——即在这两种情况下都冒犯了天主。不！你\Emph{制造}它们，使它们得以被崇拜，你\Emph{就}崇拜了它们；你崇拜，不是用某种毫无价值的香气的精神，而是用你自己的精神；不是以一头牲畜的灵魂为代价，而是以你自己的灵魂为代价。你向它们献上你的才智；你向它们奠祭你的汗水；你向它们点燃你思虑的火炬。你对于它们比司铎更是司铎，因为正是藉着\Emph{你}它们才有司铎；\Emph{你的}勤劳就是\Emph{它们的}神性。你断言你不\Emph{崇拜}你所\Emph{制造}的？啊！但那些你为他们宰杀这更肥美、更宝贵、更大的牺牲——你的救恩——的人，并不这么认为。\par

% structure: p0014 source=h2 target=section reason=relative-heading-rank; base=h2

\section{第七章 信徒对接纳制造偶像者进入教会、甚至进入圣职的悲痛}

信仰的热忱将一整天向这里陈词：哀叹一位基督徒竟从偶像进入教会；从敌对的作坊进入天主的殿；竟向天主圣父举起那作为偶像之母的手；竟用那在外受人祈祷反对天主的手向天主祈祷；竟用那赋予魔鬼躯体之手去接触主的身体。这还不够。就算从别人手中接受他们玷污之物是小事，但正是这些手将所玷污之物递给别人。制造偶像者甚至被选入教会圣秩。何其邪恶！犹太人曾将火烙加于基督；这些人则每日损伤祂的身体。啊，该被砍下的手！如今，那句话：\Quote{如果你的手使你作恶，就砍下它}\BibleRef{玛 18:8}，且看它是否仅仅是藉比喻说的。还有哪些手比那些使主的身子受绊倒的手更该被砍下呢？\par

% structure: p0016 source=h2 target=section reason=relative-heading-rank; base=h2

\section{第八章 其他服务于偶像崇拜的技艺。合法的谋生手段是充足的}

还有许多其他种类的技艺，虽然它们并不涉及\Emph{制造}偶像，却以同样的罪行\Emph{提供}那些没有它们偶像就无力的附属品。因为无论是建造还是装饰，都没有区别：如果你装饰了它的圣殿、祭台或壁龛；如果你压制了金箔，或制作了它的徽章，甚至它的房屋：这种工作所赋予的不是形状，而是权威，更为重要。如果生计的需要如此紧迫，这些技艺还有其他种类可以提供谋生手段，而不越出纪律的道路，即不制作偶像。粉刷匠既知道如何修补屋顶、涂抹灰泥、打磨水池、勾勒尖拱，以及在墙壁上浮雕许多除了肖像以外的其他装饰。画家、大理石匠、铜匠以及任何雕刻师，也都知道各自艺术的扩展，当然更容易操作。因为描绘雕像的人，给餐具柜镶嵌图案要容易得多！用椴木雕刻马尔斯神像的人，把箱子钉在一起要快得多！没有一门艺术不是相邻艺术的母亲或亲属：没有什么是独立于其相邻的。艺术的脉络如同人的贪欲一样众多。\Quote{但工资和手工报酬有差别；}因此，所需劳动也有差别。较少的工资由更频繁的收入补偿。多少墙壁需要雕像？多少为偶像建造的圣殿和神龛？但是房屋、官邸、浴室和公寓，又有多少呢？鞋和拖鞋的镀金是日常工作；但墨丘利和塞拉皮斯的镀金则不是。关于手工的收益，就此打住。奢侈和炫耀的信奉者比所有迷信都多。炫耀比起迷信更容易要求盘子和杯子。奢侈也比起仪式更多地经营花环。因此，当我们普遍劝勉人们从事那些确实不与偶像以及与偶像相称之物接触的手工时；此外，由于偶像与人共有的东西也常常是人所共有的；我们也应当提防，在我们知情的情况下，没有人从我们对偶像的服务中要求任何东西。因为如果我们做出了那种让步，又没有诉诸经常使用的补救措施，我认为我们就未能摆脱偶像崇拜的沾染，我们这些人的手（并非无意）被发现忙于照料、尊敬和服务魔鬼。\par

% structure: p0018 source=h2 target=section reason=relative-heading-rank; base=h2

\section{第九章 若干行业与偶像崇拜有关。尤其论占星术。}

我们在技艺中也观察到某些行业易受偶像崇拜的指控。至于星象学者，甚至不该提及；但如今有一个人向我们挑战，为自己坚持那行业辩护，我将略说几句。我并非指他尊崇偶像——他将偶像的名字铭刻于天上，并将天主的一切权能归于它们；因为人们推测我们受星辰不变裁决的支配，便因此认为不必寻求天主。我提出一个命题：那些背叛天主、喜爱妇女的天使，也是这种诡异艺术的发现者，并因此被天主定罪。哦，神圣的判决，其效力甚至达于大地，无知的人竟为之作证！星象学者被驱逐，正如他们的天使一样。城市和意大利禁止星象学者进入，正如天堂禁止他们的天使一样。门徒与师傅受同样的驱逐刑罚。\Quote{但贤士和星象学者从东方来了。}我们知道魔术与占星术之间的相互联盟。于是，星辰的解释者最先宣布基督的诞生，最先向祂呈献\Quote{礼物。}凭借这种关联，[必须]我想，他们使基督受惠于自己了吗？那又怎样？难道那些贤士的宗教如今也要庇护占星士吗？如今占星术果然论述基督——是基督的星象科学；而不是论及农神、战神，或其他同类死者中的任何一位，它尊崇并宣扬的乃是基督？然而，那门科学一直被允许，直到福音之时，为的是在基督诞生之后，再无人藉着天空为任何人解释他的出生。因为他们向当时还是婴儿的主献上了乳香、没药和黄金，这可以说是世俗祭献和荣耀的终结，基督将要废除它们。那么又如何？那梦——无疑出于天主的旨意——指示同一些贤士，他们应回家去，但要由另一条路，不是他们来时所走的路。这意味着：他们不应再走他们古老的道路。这并不是说黑落德不会追赶他们，事实上他并没有追赶他们；他甚至不知道他们已由\Emph{另一条路}离去，因为他本不知他们是从哪条路\Emph{来的}。正因如此，我们应当由此理解正义之路与训导。所以，那命令毋宁是：从此以后，他们应当\Emph{另走别路}。同样，另一种魔术，藉奇迹运作，甚至与梅瑟敌对，试探了天主的忍耐，直到福音之时。因为从此以后，西满·术士，刚刚成为信徒（他仍多少想着他那骗人的派系；即在他职业的奇迹中，他甚至可以藉覆手购买圣神的恩赐），便被宗徒们咒骂，并从信德中被逐出。他和另一位术士，即与色尔爵·保禄同在的那位（因为他开始反对同一位宗徒），都被处以失明的惩罚。我相信，若是有任何星象学者落入宗徒手中，也会遭受同样的命运。然而，当魔术受罚时，既然占星术是它的一个种类，那么种类自然包含在属类之中被定罪。在福音之后，你无论何处都找不到诡辩家、加色丁人、巫师、占卜者或术士，除非他们明显受罚。\BibleQuote{智者在哪里？经师在哪里？今世的辩士在哪里？天主岂不是使今世的智慧变成愚妄吗？}\BibleRef{格前 1:20}星象学者啊，你若不知道你应当成为基督徒，你就一无所知。你若知道这事，你也应当知道这一点：你应与你的那职业再无牵连，那职业本身预言他人的骤变期，并可能以自身的危险教导你。在你那体系中，你没有分，也没有产业。凡以手指或魔杖妄用天空的人，不能期望天国。\par

% structure: p0020 source=h2 target=section reason=relative-heading-rank; base=h2

\section{论教师及其困难}

此外，关于教师我们也应同样查问；不单是他们，还有所有其他文学教授者。不，相反，我们不可怀疑他们与多种拜偶像行为有密切关联：\Emph{首先}，因为他们必须宣讲外邦人的神祇，说出他们的名字、谱系、尊贵的区别，每一项；\Emph{再者}，还要遵守这些神祇的庆典和节日，正如借以计算他们收入的那些神祇的庆典。什么教师，没有七个偶像的表格，还会经常出席\OriginalTerm{克温夸特利亚节}{Quinquatria}呢？他将每个学生的首笔学费奉献给\OriginalTerm{米内瓦}{Minerva}的荣誉和名号；所以，即使他\Emph{名义上}并非被说成是\Quote{吃祭献给偶像之物}（不献给任何特定偶像），他仍被视为拜偶像者而被回避。以此为由，他招致的玷污，岂比那些名义上和实际上都公开献给偶像的行业更少呢？\OriginalTerm{米内瓦节}{Minervalia}属于米内瓦，正如\OriginalTerm{农神节}{Saturnalia}属于农神；农神节期间，连小奴隶也必须庆祝。同样，新年礼物也必须收取，\OriginalTerm{七山节}{Septimontium}也必须遵守；所有仲冬的礼物和亲族节献礼也必须索取；学校必须用花环装饰；\OriginalTerm{祭司}{flamens}的妻子们和\OriginalTerm{市政官}{aediles}献祭；学校在指定的圣日受到尊崇。同样的事发生在偶像的诞辰上；魔鬼的一切排场都被参与。谁会认为这些事情适合一位基督徒教师，除非是那些认为这些事情同样适合非教师的人？我们知道，有人会说：\Quote{如果教授文学对天主的仆人不合法，那么学习文学也同样不合法。}以及\Quote{既然文学是训练整个生命的手段，一个人如何能被训练出普通人的智力，或任何感知和行动呢？我们怎能拒绝世俗的学问，而没有它们，神圣的学问又无法进行？}那么，让我们看看文学修养的必要性；让我们思考，它部分不可接受，部分不可避免。学习文学对信徒是允许的，而非教授；因为学习与教授的原则不同。如果一个信徒教授文学，在教授时他必然推荐，在讲述时他肯定，在回忆时他见证其中所穿插的对偶像的赞美。他用这名字本身来封印那些神祇；而法律，如我们所说，禁止\Quote{宣呼众神的名号}，并将这名字归于虚妄。因此，魔鬼从他们学识的开端就建立人早期的信德。请问，那关于偶像的教导者是否犯了拜偶像之罪？但当信徒学习这些东西时，如果他已经能够理解何为拜偶像，他既不会接受也不会允许它们；如果他还不能理解，更是如此。或者，当他开始理解时，他首先应当理解他先前所学的东西，即关于天主和信德。因此他会拒绝那些东西，不会接受它们；他将如同一个人从不知情者那里故意接受毒药，却不饮下它一样安全。对他而言，必要性被当作借口，因为他没有其他学习途径。此外，不教授文学比不学习文学容易得多，正如学生不参加学校源于公开展会和学校庆典的其他污秽之事，比教师不经常参与更容易一样。\par

% structure: p0022 source=h2 target=section reason=relative-heading-rank; base=h2

\section{第十一章  贪婪与拜偶像的关系。某些行业，虽有利可图，亦当避免}

如果我们考量其余的过犯，追溯它们的来源，让我们从贪婪开始，\BibleQuote{万恶的根源，}\BibleRef{弟前 6:10}，有些人被它缠住，\BibleQuote{就在信德上翻了船。}\BibleRef{弟前 1:19}。虽然贪婪被同一位宗徒称为\OriginalTerm{拜偶像}{idolatry}。其次，\Emph{进而言之}，谎言——贪婪的仆役（关于虚誓我保持沉默，因为连发誓也不合法）——\Emph{贸易}适于天主的仆人吗？但是，撇开贪婪，获取的动机是什么？当获取的动机停止时，交易就没有必要了。暂且承认在商业中有某种义德，免除了防范贪婪和谎言的责任；我认为那种关乎偶像的灵魂和精神、喂养一切魔鬼的贸易，当属拜偶像之罪。岂不更是\Emph{首要的}拜偶像吗？如果同样的货物——我是指乳香和所有其他外国产品——被用作祭献给偶像的牺牲，同时也用于人类的药膏，对我们\Emph{基督徒}而言，还用于葬礼的慰藉，那就让他们自己斟酌吧。无论如何，当偶像的排场、偶像的司祭、偶像的祭献是通过危险、损失、不便、思虑、奔走或贸易来提供时，你被证明是什么？难道不是偶像的代理人吗？不要让任何人争辩说，这样就能对\Emph{所有}贸易提出例外。所有较严重的过犯都按危险的程度扩大警醒勤奋的范围；好使我们不仅远离过犯，而且远离过犯赖以存在的手段。因为虽然过犯由他人所做，但如果是\Emph{借我的手}，就没有区别。我绝不该成为他人的\Emph{必要}，当他做对我而言不合法的事时。因此我应当明白，我必須小心，以免我禁止做的事被\Emph{借我的手}做出来。简而言之，在另一个并不更轻的罪过上，我注意到同样的预判。由于我被禁止犯奸淫，我不会为此目的向他人提供帮助或纵容；由于我使自己肉体远离妓院，我承认我不能从事拉皮条的生意，或为邻居的利益保留那种场所。同样，禁止杀人表明角斗士教练也被排除在教会之外；没有人不成为他所供人去做之事的手段。看，这里有一个更相近的预判：如果一个公共祭牲的供应者归信，你会允许他永久从事那个贸易吗？或者，如果一个已经是信徒的人从事了那种生意，你会认为他应该留在教会里吗？不，我认为；除非有人也在乳香贩的问题上装糊涂。事实上，\Emph{血}的买卖属于一些人，\Emph{香料}的买卖属于另一些人。如果，在偶像存在于世上之前，偶像崇拜，至今尚未成形，曾通过这些货物进行；如果，即使是现在，偶像崇拜的工作多半在没有偶像的情况下，通过焚烧香料来实现；那么乳香贩甚至对魔鬼更有用，因为偶像崇拜没有偶像比没有乳香贩的货物更容易进行。让我们彻底审问信德本身的良心。一个基督徒乳香贩，如果他经过庙宇，将以什么口唾弃并吹灭那些冒着烟的祭台，这些祭台他自己曾供应？他有什么一致性去驱魔他的养子，他把自己家作为仓库提供给他们？的确，如果他驱逐了一个魔鬼，让他\Emph{自己}不要因自己的信德自夸，因为他并没有驱逐一个\Emph{敌人}；他应该容易得到他每日喂养的那一位的俯允。因此，没有任何工艺、职业或贸易，只要服务于装备或制作偶像，就能免于拜偶像的罪名；除非我们把拜偶像解释成完全不同于侍奉偶像的事。\par

% structure: p0024 source=h2 target=section reason=relative-heading-rank; base=h2

\section{第十二章 对\Quote{我该如何生活？}这一恳求的进一步回答}

我们徒然以人类维生之必需来自我安慰——若在信德被封印之后，我们说：\Quote{我没有生存的手段？}因为现在我要更充分地回应当前那唐突的主张。它提出得\Emph{太迟了}。因为效法那位最谨慎的建筑者，他先计算工程的费用以及自己的财力，免得开始之后，后来因发现自己耗尽而脸红，那本应在\Emph{之前}就考虑。但即便现在，你们有主的言语作为榜样，除去你们所有的借口。因为你说了什么？\Quote{我将陷入匮乏。}但主称匮乏者为\Quote{有福的。}\BibleRef{路 6:20}\Quote{我将没有食物。}但祂说：\Quote{不要思虑}\Quote{食物；}并以百合花作为衣着的榜样。\Quote{我的工作曾是我的生计。}不，但\Quote{一切都要变卖，分给穷人。}\Quote{但必须为子女和后代预备。}\Quote{手扶着犁向后看的，不配}做工作。\Quote{但我曾立约。}\Quote{没有人能侍奉两个主人。}你若愿意作主的门徒，就必须\Quote{背起你的十字架，跟随主：}\Emph{你的十字架}；那就是你自己的\Emph{困境}和\Emph{折磨}，或仅仅你的\Emph{身体}，它如同一个\Emph{十字架}。父母、妻子、子女，都须为天主的缘故撇下。你还在为子女和父母而犹豫于手艺、贸易、以及类似的职业吗？甚至在那里已向我们显明：既\Quote{亲爱的抵押品，}又手工活和贸易，都须为主的缘故完全撇下；当雅各伯和若望被主呼召时，他们完全撇下了父亲和船；当玛窦从税关被唤起来；甚至埋葬父亲对信仰来说也是过于迟缓的事。\BibleRef{路 9:59}主所拣选的人中，没有一人说：\Quote{我没有生存的手段。}信德不怕饥荒。它同样知道，为天主的缘故，饥饿与任何一种死亡都同样应被藐视。它已学会不重视\Emph{生命}；何况\Emph{食物}呢？〔你问〕\Quote{有多少人成就了这些条件？}但在人不可能的，在天主却容易。然而，让我们这样从天主的温和与仁慈中得到安慰：不要放纵我们的\Quote{必需}以至于接近偶像崇拜，而要远避它的一切气息，如同瘟疫一般。（并且）不仅在前述情形中，而且在整个人类迷信的系列中：无论是归于其神明，还是归于死者，或归于君王，如同属于同一不洁之灵，有时通过祭献和司祭职，有时通过表演和类似之事，有时通过圣日。\par

% structure: p0026 source=h2 target=section reason=relative-heading-rank; base=h2

\section{第13章 论与偶像崇拜有关的日期遵守}

但是，为何谈论祭献和司祭职呢？此外，关于表演和那类娱乐，我们已经用了独自一卷。在此处必须处理节日和其他特殊庆典的问题，这些有时我们出于放纵，有时出于怯懦，而违背共同的信德和纪律。确实，我要争论的第一个点是：天主的仆人是否应当与那些外邦人本身在同类之事上，无论穿着、饮食，或任何其他他们的欢乐之事分享？\Quote{与喜乐的人一同喜乐，与哀哭的人一同哀哭}\BibleRef{罗 12:15}，这是宗徒在劝勉同心合意时关于\Emph{弟兄们}所说的。但为\Emph{这些}目的，\Quote{光明与黑暗之间毫无相通}，或生与死；否则我们废除所写的：\Quote{世界将要喜乐，你们却要忧愁。}若我们与世同乐，就有理由害怕我们也将与世同忧。但当世乐时，让我们忧愁；当世后来忧愁时，我们将会喜乐。同样，阴府中的拉匝禄（在亚巴郎怀中得安慰）和那富翁（另一方面，处于火焰的折磨中）通过相应的报应，补偿了他们恶与善的交替变迁。有些赠礼日，对某些人调整荣誉的要求，对另一些人调整薪金的债务。\Quote{那么，}你却说，\Quote{我将收回我的，或偿还别人的。}若人们因迷信而为自己奉献了这一习俗，为何你，既已远离他们的一切虚荣，却参与奉献给偶像的庆典？好像对你也有某种关于某一天的规条，除了遵守某个特定日子之外，阻止你支付或收取你欠人的或人欠你的？给我你希望被对待的方式。因为为何你要躲藏，当你因邻人的无知而玷污自己的良心？若你并非不被人知道是基督徒，你受诱惑，并且你如同不是基督徒行事，违背邻人的良心；但若你也伪装，你是那诱惑的奴隶。无论如何，无论后者或前者方式，你有\Quote{以天主为耻}的罪。但\Quote{凡在人前以我为耻的，我在我天上的父前也必以他为耻，}祂说，\Quote{在我天上的父面前。}\par

% structure: p0028 source=h2 target=section reason=relative-heading-rank; base=h2

\section{第14章 论亵渎。圣保禄的一句话}

但是，迄今大多数（基督徒）心中已经产生这样一种信念：认为他们若有时跟外邦人行事，是可以宽恕的，因为害怕\Quote{那名被亵渎}。然而，在我看来，必须千方百计躲避的亵渎是这样的：如果我们中的任何人因欺诈、伤害、侮辱或其他值得抱怨的事，有正当理由导致外邦人亵渎，在其中\Quote{那名}理应受到抨击，以致主也应理当发怒。否则，如果关于一切的亵渎都说过：\Quote{因你们的缘故，我的名被亵渎}，那么我们立刻都灭亡了；因为整个竞技场，尽管我们毫无过失，却用邪恶的票选攻击\Quote{那名}。让我们停止（作基督徒），那它就不会被亵渎！相反，只要我们存在，就让它被亵渎：是在遵守纪律，而非逾越纪律之时；是在我们被认可时，而非被摈弃时。哦，这块近殉道的亵渎，现在它\Emph{证明}我是基督徒，然而正因如此，它又\Emph{憎恶}我！对良好纪律的咒骂乃是那名的祝福。他说：\Quote{如果我愿意讨人的喜欢，我就不是基督的仆人了。}但同一位宗徒在别处却吩咐我们要小心讨所有人喜欢：\Quote{像我一样，在一切事上讨众人喜欢。}\BibleRef{格前 10:32}无疑，他那时是通过庆祝\OriginalTerm{农神节}{Saturnalia}和\OriginalTerm{元旦}{New-year's day}来讨他们喜欢吗？［是这样吗？］还是通过节制和忍耐？通过严肃、仁慈、正直？同样，当他说：\Quote{对一切人，我就成为一切，为的是要赢得一切人，}\BibleRef{格前 9:22}他是指\Quote{对拜偶像的人成为拜偶像者？}\Quote{对外邦人成为外邦人？}\Quote{对世俗的人成为世俗的？}然而，虽然他并不禁止我们与拜偶像者、奸淫者以及其他罪人来往，说：\Quote{不然，你们就得出离世界，}\BibleRef{格前 5:10}他当然没有那样放松交往的缰绳，以致既然我们有必要与罪人一同\Emph{生活}和\Emph{交往}，我们也能与他们一同\Emph{犯罪}。凡有生活的交往——宗徒允许的——就有犯罪，这是无人允许的。与外邦人同活是合法的，与他们同死则不然。让我们与所有人同活；让我们因本性的相通（而非因迷信）与他们一同喜乐。我们在灵魂上平等，不在纪律上；我们是世界的共同拥有者，不是错误的共同拥有者。但如果我们与陌生人在这类事情上没有共融的权利，那么在弟兄中庆祝它们是多么更邪恶啊！谁能维持或辩护这一点？圣神责备犹太人的庆节。祂说：\Quote{你们的安息日、月朔和节期，我的心恨恶。}而我们——安息日对我们来说是陌生的，月朔和从前为天主所爱的节期也是——却频繁庆祝\OriginalTerm{农神节}{Saturnalia}、\OriginalTerm{元旦}{New-year's}、\OriginalTerm{仲冬节}{Midwinter's festivals}和\OriginalTerm{主妇节}{Matronalia}——礼物来来往往，新年赠礼，游戏嘈杂，宴席喧闹！哦，外邦人对自己的教派有更好的忠诚，他们不会要求基督徒的庄严礼仪！他们即使知道主日和\OriginalTerm{五旬节}{Pentecost}，也不会与我们共享；因为他们害怕被人看作基督徒。\Emph{我们}却不担心被人看作\Emph{外邦人}！若许给肉身任何放纵，你已有了。我不说你们自己的日子，而是更多；因为对外邦人来说，每个节期一年只一次；而\Emph{你们}每八天就有一个节期。请列举外邦人的个别庆典，将它们排成一列，它们也无法凑成一个\OriginalTerm{五旬节}{Pentecost}。\par

% structure: p0030 source=h2 target=section reason=relative-heading-rank; base=h2

\section{第十五章。关于皇帝、胜利等节庆。\Person{三青年}和\Person{达尼尔}的榜样。}

祂说：\Quote{让你们的善行发光}；\BibleRef{玛 5:16}但如今我们的店铺和门户却发光！如今你会找到更多外邦人的门没有灯和桂冠，反而不见基督徒的门。这种情况对这种仪式形式又意味着什么？如果这是对偶像的尊崇，那么无疑对偶像的尊崇就是偶像崇拜。如果是为人的缘故，让我们再思考一下：所有偶像崇拜都是为人的缘故；让我们再思考：所有偶像崇拜都是对人进行的敬拜，因为即使在敬拜者中间也普遍承认，从前万国的神本身也是人；所以，将这种迷信的敬意献给前代的人或当代的人并无区别。偶像崇拜被定罪，不是因为被敬拜的对象，而是因为这些属于魔鬼的礼节。\Quote{凯撒的当归给凯撒}。祂将这句话与\Quote{天主的当归给天主}并列，就足够了。那么，什么是凯撒的？就是当时所咨询的事：是否应纳人头税给凯撒？因此，主也要求将钱币指给他看，并询问上面的像是谁的；当祂听说那是凯撒的，便说：\Quote{凯撒的当归给凯撒，天主的当归给天主}；也就是说，钱币上凯撒的像归凯撒，人身上天主的像归天主；这样，你把钱给凯撒，把自己给天主。否则，如果一切都是凯撒的，什么是天主的？那么，你可能会说：\Quote{我门前的灯和门柱上的桂冠是对天主的尊崇？}\Emph{它们在那里}，当然不是因为它们是对天主的尊崇，而是对那藉这种仪式礼仪代替天主位置者的尊崇，就显明而言，但暗中的宗教行为却是属于魔鬼的。因为我们应当确知，若有任何人因对世间文献的无知而忽略了这一点，那就是在罗马人中间也有门神：按铰链得名的\OriginalTerm{卡耳德亚}{Cardea}（铰链女神），按门得名的\OriginalTerm{福耳库卢斯}{Forculus}（门神），按门槛得名的\OriginalTerm{利门提努斯}{Limentinus}（门槛神），以及\OriginalTerm{雅努斯}{Janus}本人（门神）；当然我们知道，尽管这些名字是空洞和虚构的，但当它们被引入迷信时，魔鬼和一切不洁之灵便通过奉献的关联抓住它们。否则魔鬼没有个体名字，但他们在那里找到一个名字，就像找到一个记号一样。在希腊人中间，我们也读到\OriginalTerm{阿波罗·提瑞俄斯}{Apollo Thyræus}，即门之神，以及\OriginalTerm{安忒利或安忒利诸魔}{Antelii or Anthelii}，作为入口的守护者。因此，圣神从起初就预见这些，通过最古老的先知厄诺克预先歌唱，说连入口也会成为迷信之用。因为我们看到在浴场中其他入口也被崇拜。但如果有受崇拜的入口，那么灯和桂冠就属它们。你无论对入口做了什么，都是对偶像做了什么。在此，我也呼吁天主权威作证；因为隐瞒任何向某人显示的事（当然是为众人的缘故）是不安全的。我知道有一位兄弟在当夜通过异象受到严厉责罚，因为当突然宣布公共庆祝时，他的仆人用花环装饰了他的大门。而他自己并没有装饰，也没有命令装饰；因为他之前出门，回来后责备了这件事。在这种事上，我们与天主如此严格地被审视，甚至关乎我们家庭的纪律。因此，关于对君王或皇帝的尊崇，我们有足够的规范，即我们应当完全服从，依据宗徒的训令：\Quote{服从官员、君侯和掌权者}；\BibleRef{铎 3:1}但限于纪律的范围，只要我们保持自己与偶像崇拜分离。因为正是为此，那三位兄弟的榜样走在我们前面，他们在其他方面服从拿步高王，却坚定地拒绝敬拜他的像，\EditorialNote{达 2-3}证明任何被抬高超过人类尊荣的限度、达到神圣崇高之相似的事物，就是偶像崇拜。同样，达尼尔在其他方面顺从达理阿，但在不危及宗教信仰的范围内，他一直遵守本分；\EditorialNote{达 6}因为他为避免那危险，不怕王狮，正如他们不怕王火。因此，让那些没有光的人每日点灯；让那些地狱之火临近的人把注定要燃烧的桂冠贴在门柱上：黑暗的见证和刑罚的预兆适合他们。\Emph{你}是世界的光，是常青的树。如果你已弃绝庙宇，就不要把自己的门变成庙宇。我说得太少了。如果你已弃绝妓院，就不要用新妓院的外表装饰你自己的房子。\par

% structure: p0032 source=h2 target=section reason=relative-heading-rank; base=h2

\section{第十六章 论私人节庆}

然而，关于私人及社交典礼的礼仪——如白色托加袍、订婚、结婚、命名仪式——我认为不必担心其中掺杂的敬拜偶像的气味。因为应当考虑这些礼仪所针对的原因。上述礼仪在我看来本身是洁净的，因为男子的服饰、订婚戒指或婚姻的结合，并非源自对任何偶像的尊崇。简而言之，我没有发现任何服饰被天主咒诅，除了男子穿女子服装：因为祂说：\Quote{被咒骂的}，祂说，\Quote{是每一个穿女人衣服的男人。}然而，托加袍无论从\Emph{名称}上还是从\Emph{用途}上都是男子的服饰。天主并不禁止举行婚礼，也不禁止取名。\Quote{但这些场合有专属的祭献。}若我被邀请，且典礼的名称不是\Quote{参与祭献}，那么我履行对朋友的服务。但愿它真是\Quote{为朋友服务}，使我们能避免看见那些我们不该做的事。但由于那恶者已以敬拜偶像如此环绕世界，参与一些只关乎服侍人而非服侍偶像的典礼，对我们而言是合法的。显然，若我被邀请参与司祭职务和祭献，我不会去，因为那是专属偶像的服务；但我也不会就此类事情提供建议、费用或任何其他帮助。若我因祭献而被邀请并站在一旁，我将成为敬拜偶像的共犯；若任何其他原因使我和祭献者在一起，我仅是祭献的旁观者。\par

% structure: p0034 source=h2 target=section reason=relative-heading-rank; base=h2

\section{第十七章 仆役与其他官员的情形。基督徒可担任何种职务}

但信主的仆役或子女该如何呢？同样，当官员们侍奉他们的主人、恩主或上级行祭献时，又该如何？若有人把酒递给献祭者，甚至若他用任何与祭献相关或属于祭献的言语协助了献祭者，他就会被视为敬拜偶像的仆役。铭记这条规则，我们可以效法先祖和其他先辈的榜样，甚至向\Quote{官长和掌权者}提供服务，他们服从敬拜偶像的君王，直至敬拜偶像的边界。最近因而产生了一项争论：天主的仆役是否应担任任何尊位或权力，如果他能够——或藉某种特殊恩宠，或藉灵巧——使自己免受各种敬拜偶像的玷污；仿效若瑟和达尼尔的榜样，他们远离敬拜偶像，在埃及全境或巴比伦的首府穿着总督的制服和紫袍，掌管尊位和权力。因此，我们承认，任何人都有可能在任何职位上，仅以职位的\Emph{名义}行事：既不献祭，也不将自己的权威用于祭献；不出租牺牲；不把神庙的管理指派给别人；不征收其贡品；不自费或公费举办或主持表演；不为任何庆典发布通告或谕令；甚至不宣誓；此外（属于\Emph{权力}范畴的）：不判决任何人的生命或品行，因为尚可容忍他对\Emph{金钱}作判决；不定罪也不预判；不捆绑任何人，不监禁或拷打任何人——若这一切确有可能。\par

% structure: p0036 source=h2 target=section reason=relative-heading-rank; base=h2

\section{第十八章 与敬拜偶像相关的服装}

但现在我们只须论及职位的服饰和装备。每个人都有其适宜的衣着，既用于日常，也用于职务和尊严。因此，那著名的紫袍和金饰作为颈上的装饰，在埃及人和巴比伦人当中，曾是尊严的标志，正如如今有镶边、条纹或绣棕榈叶的长袍，以及省司铎的金冠；但条件并不相同。因为过去它们只以\Emph{荣誉}之名授予那些配得国王亲密友谊的人（正因如此，那些人常被称为国王的\Quote{紫衣人}，正如在我们当中，有些人因白袍而被称为\Quote{候选人}）；但\Emph{并非}意味着那衣着也应与\Emph{司铎职}或\Emph{任何偶像仪式}绑定。因为如果\Emph{那样}的话，当然，如此圣洁和坚定的人会立即拒绝那些被玷污的服饰；而且也会立即显明达尼尔并非热衷崇拜偶像，也未敬拜贝耳或那条龙，这些在很久以后\Emph{确实}显明了。因此，那紫袍是单纯的，当时在蛮族中并非尊严的标志，而是\Emph{贵族}的标志。因为正如曾是奴隶的若瑟和因被掳而改变身份的达尼尔，通过蛮族贵族的服饰获得了巴比伦和埃及的自由；同样，在我们信徒中，若有需要，镶边袍可适当让给男孩，披肩给女孩，作为出身而非权势的标志；种族而非职务的标志；等级而非迷信的标志。但那紫袍或其他尊严与权力的标志，从一开始就奉献给附在尊严和权力上的偶像崇拜，带有其自身亵渎的污点；因为，此外，镶边和条纹的长袍以及宽条纹的袍子，甚至被用在偶像本身上；\OriginalTerm{束棒}{fasces}和棍棒也放在它们前面；这理所当然，因为魔鬼是这世界的执政者：它们拿着\Emph{束棒}和紫袍，属于同一团体的标志。那么，如果你使用衣着却不执行其职能，你会有什么进展？在不洁的事物中，无人能显现为洁净。如果你穿上本身被玷污的束腰外衣，它也许不会通过你而被玷污；但你通过它，将无法洁净。现在，那些以\Quote{若瑟}和\Quote{达尼尔}为论据的人要知道，新旧事物、粗俗与精致、开始与发展、奴役与自由，并不总是可比的。因为他们，即便按他们的处境，也是奴隶；但你，不为任何人所奴役，就你仅为基督的奴隶而言，祂也已将你从世界的囚禁中释放出来，你将承担按照主的榜样行事的责任。那主行走在谦卑和隐晦中，没有固定的住所：因为\Quote{人子，}祂说，\Quote{没有枕头的地方；}衣着朴素，否则祂就不会说：\Quote{看，穿华丽衣服的人是在王宫里；}简言之，容貌和外表无荣耀，正如依撒意亚也曾预言过。\BibleRef{依 53:2}如果祂甚至对自己的门徒也没有行使权力，反而向他们提供仆役般的服务；如果，简言之，尽管祂知道自己的王国，却退避而拒绝被立为王，\BibleRef{若 6:15}祂就以最完全的方式给祂自己的人一个榜样，要他们冷漠地对待一切尊严和权力的骄傲和服饰。因为\Emph{如果它们要被使用}，谁比天主之子更会使用它们呢？何种、多少\Emph{束棒}会伴随着祂？何种紫袍会在祂肩上绽放？何种金色会在祂头上闪耀，如果祂没有断定世界的荣耀对祂自己和祂的人都是不相干的？因此祂不愿意接受的，祂就拒绝了；祂拒绝的，祂就谴责了；祂谴责的，祂就视为魔鬼排场的一部分。因为祂不会谴责，除非那些不属于祂的；而不属于天主的，只能是魔鬼的。如果你已弃绝\Quote{魔鬼的排场}，要知道凡你所触摸的那里的一切都是偶像崇拜。让这个事实也提醒你：这世界的一切权力和尊严不仅与天主相异，而且是祂的仇敌；藉着它们，对天主仆人的惩罚已被决定；藉着它们，为不敬者准备的刑罚也被忽视了。但\Quote{你的出身和财富在抵制偶像崇拜时让你困扰。}为避开它，补救办法不会缺乏；因为即使缺乏，仍存留着那一个，藉此你将成为更幸福的官长，不是在地上，而是在天上。\par

% structure: p0038 source=h2 target=section reason=relative-heading-rank; base=h2

\section{第十九章  论军事服役}

在上一节中，决定似乎也已涉及军事服役，它介于尊严和权力之间。但现在要探究的是，信徒是否可以投身军事服役，以及军人是否可以接受信仰，即使是士兵或每个下级军阶，他们无需参与祭祀或死刑。神圣的圣事与人的圣事、基督的旗帜与魔鬼的旗帜、光明的营地与黑暗的营地之间没有一致。一个灵魂不能同时向两个\Emph{主人}——天主和凯撒——负责。然而梅瑟曾执杖，亚郎佩带纽扣，若翰（施洗者）束皮带，农的儿子若苏厄领队行军；百姓也打过仗：如果你喜欢玩弄这个话题。但\Emph{一个基督徒}如何作战？他如何在和平时期服役，却没有一把主所除去的剑？因为虽然士兵曾来到若翰那里，并接受了他们规则的标准；虽然，同样，一位百夫长曾相信；\Emph{然而}主后来在解除伯多禄的武装时，解开了每一个士兵的带子。在我们当中，任何被指派给非法行为的服饰都是不合法的。\par

% structure: p0040 source=h2 target=section reason=relative-heading-rank; base=h2

\section{第二十章  论言语中的偶像崇拜}

然而，由于符合神圣规范的行为不仅因行为而受危，也因言语而受危（因为正如经上记载：\Quote{看，这人和他的作为；}同样，\BibleQuote{凭你的话，你将被称义}\BibleRef{玛 12:37}），我们应当记住，即使在\Emph{言语}上，也必须防范偶像崇拜的侵入，这侵入或是出于习俗的缺陷，或是出于怯懦。律法禁止提说外邦人的神，当然不是指我们不可说出他们的名字——日常交往迫使我们要说出这些名字：因为常常必须说：\Quote{你在厄斯库拉庇乌斯庙里找到他；}以及\Quote{我住在伊西斯街；}以及\Quote{他被立为朱庇特的司铎；}还有许多类似的情况，因为甚至\Emph{人}也经常被赋予这类名字。如果我按他的本名称一个人为萨图尔努斯，我并不敬拜萨图尔努斯；我敬拜他不比称一个人为马尔库斯时敬拜马尔库斯更多。但经上说：\BibleQuote{不可提说其他神的名字，也不可从你口中听见。}\BibleRef{出 23:13}它所给予的诫命是，我们不可称它们为神。因为在律法的第一部分，祂也说：\Quote{你不可，}祂说，\Quote{妄称主你天主的名，}\BibleRef{出 20:7}即用于偶像身上。因此，凡以天主之名尊崇偶像者，已陷入偶像崇拜。但如果我说到它们时称它们为神，必须加上一些东西以示我自己并非在称它们为神。因为连圣经也提到\Quote{神，}但加上\Quote{他们的，}即\Quote{外邦人的；}正如达味在提到\Quote{神}时所说的：\Quote{外邦人的神都是邪魔。}但我这样说是为后续的观察奠定基础。不过，习俗的缺陷在于说：\Quote{凭赫丘利，愿信义之神帮助我；}而在习俗之上又加上了一些人的\Emph{无知}，他们不知道那是凭赫丘利起誓。此外，以你已弃绝的神之名起誓，岂不成了信德与偶像崇拜的勾结？因为谁不以他所起誓之名为尊呢？\par

% structure: p0042 source=h2 target=section reason=relative-heading-rank; base=h2

\section{第21章 论对外邦人言词的默许}

但这乃是怯懦的标记：当别人以他的神之名，通过起誓或其他形式的证词来约束你时，你因害怕暴露而保持沉默。因为同样，你因沉默而肯定了他们的威严，基于这威严，你似乎被约束。这有什么区别呢？你称外邦人的神为神，或听他们被称为神，难道不是肯定他们？你以偶像起誓，或当别人命令你起誓时你默许，难道不是一样？我们为何不识别撒殚的诡计？他旨在通过我们的口不能成就的事，通过他仆人的口来成就，借着我们的耳朵将偶像崇拜引入我们。无论如何，无论命令者是谁，他是在友好或敌对的联合中约束你。如果是敌对的，你已被挑战去战斗，并且知道你必须作战。如果是友好的，你将何等安全地将你的承诺转移给主，以便解除那人的约束——那恶者正通过他试图使你依附于偶像的尊荣，即偶像崇拜！所有这类忍受都是偶像崇拜。你尊敬那些当被强加为权威时你曾表示尊敬的对象。我知道一人（愿主宽恕他！），当在诉讼中有人公开对他说：\Quote{愿朱庇特对你发怒，}他回答说：\Quote{相反，愿他对\Emph{你}发怒。}一个相信朱庇特是神的外邦人会怎么做呢？即使他没有以朱庇特之名回咒（或其他类似的话），但\Emph{仅仅}通过回咒，他就确认了朱庇特的神性，显示出因朱庇特之名的咒骂而恼怒。因为如果你知道某人不存在，你为何会在以他之名的咒骂下愤怒呢？如果你发怒，你立刻肯定了他的存在，而你恐惧的宣认便成了偶像崇拜的行为。更何况，当你以朱庇特之名回咒时，你尊崇朱庇特的方式与激怒你的人相同！但信徒应当在这种情况下发笑，而非发怒；不仅如此，按照诫命，甚至不以天主之名回咒，而要亲切地以天主之名祝福，以便你既能摧毁偶像，又能宣讲天主，并成就纪律。\par

% structure: p0044 source=h2 target=section reason=relative-heading-rank; base=h2

\section{第22章 论接受以偶像之名祝福}

同样，一个已加入基督的人不会容忍以外邦人之神的名被祝福，而总会拒绝这不洁的祝福，并通过将其转向天主来为自己净化。以外邦人之神的名被祝福，就是以天主之名被咒骂。如果我施舍了财物，或行了其他善事，而受惠者祈求他的神或殖民地的护守神对我垂怜，我的奉献或行为便立即成了对偶像的尊荣，因他以偶像之名回报我的祝福。但他为何不知道我是为天主而行呢？好使天主更受光荣，而非使魔鬼因我为天主所做的而被尊崇。如果天主看见我为祂而行，祂也同样看见我不愿意表明我是为祂而行，并在某种程度上将祂的诫命献给了偶像。许多人说：\Quote{没有人应当暴露自己；}但我认为也不应当\Emph{否认}自己。因为凡在任何事上伪装而被视为外邦人者，就是在否认；而当然，一切否认都是偶像崇拜，正如一切偶像崇拜都是否认，无论在行为上或言语上。\par

% structure: p0046 source=h2 target=section reason=relative-heading-rank; base=h2

\section{第23章 以偶像之名所立之书面契约 默示同意}

但有一类行为，在言行上双倍锋利，两方面都有害处，尽管它谄媚你，仿佛各面都无害；因为它不被视作\Emph{行为}，因为它不被当作\Emph{言语}来抓住。在向异教徒抵押借款时，\Emph{基督徒}立下誓言担保，并否认自己\Emph{如此做了}。当然，起诉的时间、审判庭的地点、主审法官的身份，决定了他们知道自己\Emph{如此做了}。基督规定不可起誓。\Quote{我写了，}债务人说，\Quote{但我什么也没说。是舌头，不是所写的文字，在杀人。}在此我呼号自然和良心为我的证人：自然，因为即使在口授时\Emph{舌头}保持静止和沉默，手也不能写出\Emph{灵魂}未曾口授的任何东西；尽管连对舌头本身，灵魂也可能口授了它自己构思的东西，或由他人传达的东西。现在，以免有人说：\Quote{是另一个人在口授，}我在此诉诸良心：他人所口授的，灵魂是否接纳并传递给手，无论舌头是伴随还是不动。主说过，过错是在心思和良心中犯下的。如果私欲或恶意上升到一个人心中，祂说这就是被视为行为。因此你立了担保；这显然\Quote{上升到你的心中}，你不能争辩说你不知道或不愿意；因为当你立担保时，你知道\Emph{你做了}；当你知道了，当然你是愿意的：你既在行为上也在思想上做了；你也不能用较轻的指控排除较重的，以致说它显然是虚假的，因为你立了担保却不实际履行。\Quote{但我没有否认，因为我并没有起誓。}\Emph{但你起了誓}，因为即使你并未如此做，只要你\Emph{同意}如此做，你仍可说起了誓。在\Emph{书写}的情形中，声音的沉默是无用的辩解；在\Emph{文字}的情形中，声响的沉默也是无用的。因为匝加利亚在暂时失去\Emph{声音}受罚时，与他的\Emph{心思}交谈，越过无用的舌头，借助双手从心里口授，不用嘴就宣布了他儿子的名字。因此，在他的笔下，有一只手比任何声音更响亮；在他的蜡版上，有一封信比任何嘴更可听见。要查问一个人是否\Emph{说了话}，如果他被认为是\Emph{说了话}。我们祈求主，愿那种契约的需要永不包围我们；如果\Emph{它}果真发生，愿祂赐给我们弟兄援助我们的方法，或赐给我们恒心断绝一切\Emph{这种}需要，以免那些代替我们嘴巴的否认文字，在审判之日被带出来反对我们，现在不是由证人，而是由天使以印章封印！\par

% structure: p0048 source=h2 target=section reason=relative-heading-rank; base=h2

\section{第24章 总结}

在这些崇拜偶像的礁石与海湾、浅滩与海峡之间，信德，被天主圣神的风鼓满了帆，航行着；谨慎则安全，警觉则稳妥。但那些被冲下船的人，是深不见底的汪洋，无路可游；那些搁浅的人，是无法脱险的船难；那些被吞没的人，是漩涡，无法呼吸——甚至在崇拜偶像中。它的所有波浪都令人窒息；每一个漩涡都把人吸到下阴府。没有人可以说：\Quote{谁能如此安全地防范自己？我们将不得不离开世界！}\BibleRef{格前 5:10}好像离开世界不比站在世界作为崇拜偶像者更有价值！没有比以敬畏偶像崇拜为首要恐惧更容易防范的了；与如此危险相比，任何\Quote{必要}都太微不足道了。圣神之所以在宗徒们当时商议时为我们放松了束缚和轭，\BibleRef{宗 15:1}是为了我们能自由地献身于避免偶像崇拜。这将是我们的法律，越容易使用就越要全面执行；（一条法律）专属于基督徒，藉此异教徒认出了我们并查验我们。这条法律必须放在那些接近信仰的人面前，并灌输给那些进入信仰的人；以便他们在接近时能深思熟虑；遵守它则能持守；不遵守则放弃其名号。我们将留意，如果在方舟的预象之后，教会中有乌鸦、鸢、狗和蛇。无论如何，在方舟的预象中找不到崇拜偶像者：没有动物被造来代表崇拜偶像者。愿那在方舟中没有的东西不在教会中。\par

\SourceRecord{Translated by S. Thelwall.；Ante-Nicene Fathers；Vol. 3.；Edited by Alexander Roberts, James Donaldson, and A. Cleveland Coxe.；Buffalo, NY: Christian Literature Publishing Co.,；1885.；<http://www.newadvent.org/fathers/0302.htm>.}
